\section{The library migration}
\label{sec:libraries}

The repository was written against Lean \texttt{v4.34.0-rc2} and the \Mathlib{} commit that
toolchain fixes. This section documents its migration to the current library generation and
to \TauCeti{}, and reports what the move cost.

\subsection{The target pin}
\label{sec:pin}

Two upstream constraints determine the possible targets. \TauCeti{} pins a toolchain and a
\Mathlib{} commit, and Lake refuses to build a package whose dependencies disagree about the
toolchain. At the time of writing \TauCeti{} is on Lean \texttt{v4.36.0-rc1} with \Mathlib{}
at \texttt{2cabef97}; the newest \Mathlib{} release is \texttt{v4.36.0-rc2}, whose toolchain
no \TauCeti{} release uses. We therefore chose the pin that is simultaneously the newest
\TauCeti{} and the newest \Mathlib{} that \TauCeti{} supports:
\begin{center}
\begin{tabular}{@{}lll@{}}
\toprule
component & old pin & new pin \\
\midrule
Lean & \texttt{v4.34.0-rc2} & \texttt{v4.36.0-rc1} \\
\Mathlib{} & \texttt{v4.34.0-rc2} & \texttt{2cabef97b63d} \\
\TauCeti{} & --- & \texttt{b8cbc774e9a4} \\
Comparator & \texttt{v4.34.0-rc2} & dropped (see below) \\
\bottomrule
\end{tabular}
\end{center}
The \lean{rev} fields are full commit hashes rather than branch names, so that
\lean{lake update} is deterministic; \lean{lakefile.toml} records the reason for each pin in
a comment.

\paragraph{Comparator.} Comparator is released per toolchain, and its releases stop at
\Lean{} \texttt{v4.35.0}, one toolchain below the one \TauCeti{} requires. Since Comparator is
needed only for the independent re-checking of the four headline theorems and not for their
proofs, we removed it from the dependency set. The two challenge modules and their JSON
configurations remain in the repository unchanged---they import only \Mathlib{}---so
independent checking is still possible with a Comparator installation matching whichever
toolchain is in use. This is the one place where the migration cost functionality rather
than just time, and it is recorded in \file{MIGRATION.md}.

\subsection{Practical obstacles}
\label{sec:obstacles}

The migration is not a version bump; it is a small research project in code archaeology. The
obstacles we met, in decreasing order of the number of sites they affected, were the
following.

\paragraph{\lean{MemLp} became a definition.} In the pinned \Mathlib{}, \lean{MemLp f p μ}
is \emph{defined} to be \lean{eLpNorm f p μ < ∞} rather than a structure with a
measurability field and a finiteness field. Consequently:
\begin{itemize}
\item field projections \lean{hf.1}, \lean{hf.2} must become the lemmas
      \lean{hf.aestronglyMeasurable} and \lean{hf.eLpNorm\_lt\_top};
\item anonymous-constructor terms \lean{⟨hmeas, hbound⟩} of type \lean{MemLp f p μ} must
      become the finiteness proof alone;
\item any \lean{obtain ⟨hmeas, hfin⟩ := hf} must be rewritten.
\end{itemize}

\paragraph{Measurability hypotheses moved in and out of the \texorpdfstring{$L^p$}{Lp} API.}
This is the subtle one. \lean{eLpNorm} used to be $0$ for non-a.e.-strongly-measurable
functions, which made monotonicity lemmas unconditional; it is now $\infty$ in that case,
so \lean{eLpNorm\_mono\_ae} gained an explicit measurability hypothesis and
\lean{eLpNorm\_one\_eq\_lintegral\_enorm} became a lemma \emph{with} a hypothesis. In the
other direction, \lean{eLpNorm\_add\_le} \emph{lost} its two measurability arguments and now
takes only the exponent hypothesis \lean{1 ≤ p}. A proof that says
\lean{eLpNorm\_add\_le hf.aestronglyMeasurable hg.aestronglyMeasurable hp1} must become
\lean{eLpNorm\_add\_le hp1}, while a proof that says \lean{eLpNorm\_mono\_ae h} must become
\lean{eLpNorm\_mono\_ae hf h}---possibly with a case split, because the measurability of the
dominated function is genuinely no longer automatic. One project lemma
(\lean{lpNorm\_mono\_of\_norm\_le\_ae}) acquired a proof by cases on measurability for exactly
this reason, and a hypothesis of the corresponding library statement was thereby
strengthened: a semantic change forced by the library, not a syntactic one.

\paragraph{One change is not a repair but a gap.} The most interesting failure
found by the migration is in \file{Euler/EulerProof.lean}. Five proofs there apply
\lean{eLpNorm\_le\_mul\_eLpNorm\_of\_ae\_le\_mul} to an \emph{iterated field derivative}, with no
measurability hypothesis in hand, because under the old convention \lean{eLpNorm} was $0$ on
functions that are not a.e.\ strongly measurable and the inequality was therefore free. Under
the pinned library \lean{eLpNorm} is $\infty$ there, so the lemma needs measurability of that
derivative, and the development has no lemma producing it: the smoothness it proves is
smoothness of the \emph{lifted local field}, not measurability of the derivative-valued map on
the domain. The proofs are not syntactically broken; they are correct only under an
assumption the library no longer makes, and closing them means proving a measurability lemma
that the development never needed. This is the clearest instance in the whole migration of a
library change that alters what has to be proved rather than how it is written, and it is why
the residual error count in \file{EulerProof} should not be read as a measure of remaining
effort.

\paragraph{Explicit exponent hypotheses became a structure.}
\lean{eLpNorm'\_le\_eLpNorm'\_mul\_eLpNorm'} (the interpolation inequality for
\lean{eLpNorm'}) previously took the three exponent facts as separate arguments; it now takes
a single \lean{Real.HolderTriple p q r}, and its two measurability arguments moved from the
front to the middle of the argument list. Interpolation sites must be rewritten to build the
structure:
\begin{lstlisting}
have hpqr : Real.HolderTriple (p.toReal / a) (q.toReal / b) r.toReal :=
  { inv_add_inv_eq_inv := by simpa only [one_div] using hab.symm
    left_pos := div_pos hp hpt ha
    right_pos := div_pos hq hqt hb }
simpa using eLpNorm'_le_eLpNorm'_mul_eLpNorm'
  (fun x y : ℝ => x * y) 1
  (hf.norm.aemeasurable.pow_const a).aestronglyMeasurable
  (hg.norm.aemeasurable.pow_const b).aestronglyMeasurable
  (Filter.Eventually.of_forall fun x => by simp [nnnorm_mul])
  hpqr
\end{lstlisting}

\paragraph{Hidden arguments became explicit and vice versa.}
\lean{contDiff\_apply} now takes the differentiability order implicitly, so every application
in a fixed-order context must say \lean{(n := ∞)}; \lean{norm\_iteratedFDeriv\_fderiv}
changed from ``$\|D^n(\mathrm{d}f)\| = \|D^{n+1}f\|$'' to a form in which the dependent
function is implicit, so a two-argument application must be replaced by an explicitly typed
\lean{have}; \lean{Finset.prod\_le\_prod} now demands a \lean{MulLeftMono} instance that
$\mathbb{R}$ does not have, so a monotonicity proof for a product of factors of the form $b+k$
must be restated so that positivity of the factors is available.

\paragraph{The module system.}
\Lean{} \texttt{v4.36} introduces the module system: a file opened with \lean{module} may
\lean{public import} and hides its non-public declarations. \Mathlib{} and \TauCeti{} have
adopted it, and Lean emits a warning for every import of a module file from a file that has
not declared itself a module:
\begin{lstlisting}
warning: NavierStokes/R3/ComparisonGronwall.lean: imports
  `Mathlib.Analysis.Calculus.Deriv.MeanValue` from package `mathlib`, which is designed for
  use with the module system; consider adding `module` to the start of this file
\end{lstlisting}
The warnings are harmless, but they are also loud: at this scale they run to several per
file. Adding a \lean{module} header to every file makes them disappear and puts the
repository on the same footing as its dependencies; a \lean{module} file cannot import a
non-\lean{module} file, so the change is all-or-nothing.

\subsection{\TauCeti{}: what it supplies}
\label{sec:tauceti}

\TauCeti{} is a downstream analysis library organised along the same lines as \Mathlib{}
itself: 2{,}793 files under \file{TauCeti/Analysis/}, of which the areas relevant here are
Sobolev theory (142 files: weak derivatives, Bessel potentials, embeddings, mollification,
Poincaré, Rellich--Kondrachov, trace theory on Lipschitz domains), PDE (161 files:
fundamental solutions, Green functions, the heat kernel with its semigroup and Duhamel
potential, Caccioppoli and Harnack estimates), distributions (56 files: Schwartz space,
Littlewood--Paley, tempered distributions, Fourier multipliers), Fourier theory (42 files)
and parametric integration. It is built on \Mathlib{} and uses the module system throughout.

For this repository the interesting entry points are these.

\begin{center}\small
\begin{longtable}{@{}p{0.36\textwidth}p{0.58\textwidth}@{}}
\toprule
\TauCeti{} module & relevance \\
\midrule
\file{Analysis/Calculus/ParametricIntegral} & $C^n$ regularity of compact-parameter
integrals, iterated derivatives under an interval integral with individually bounded
derivatives; replaces the project's \file{SmoothParameterIntegral} and
\file{CompactParameterIntegral} \\
\file{Analysis/Sobolev/Mollification} & the mollification identity
$D_v(u\star\rho) = u'\star\rho$ for weak derivatives, followed by density \\
\file{Analysis/Sobolev/Embedding}, \file{Morrey}, \file{LipschitzDomain/*} & Sobolev
embeddings, trace and extension theorems on Lipschitz domains \\
\file{Analysis/PDE/HeatKernel/*} & heat kernel, semigroup, smoothing, Duhamel potential
--- the analytic core of the exterior swirl field \\
\file{Analysis/Distribution/SchwartzSpace/LittlewoodPaley} & frequency localisation at the
level the dyadic charts use informally \\
\file{Analysis/Fourier/MihlinHormander} & Fourier-multiplier bounds for the pressure
recovery \\
\bottomrule
\end{longtable}
\end{center}

\subsection{What can actually be replaced}
\label{sec:replace}

It is easy to overstate this. The project's analysis is not organised the way a textbook
library is organised: its Sobolev spaces are word lattices over a cylinder designed to be
complete \emph{as defined}, its cutoffs are built with explicit all-order bounds tied to
reserved logarithmic patches, and its heat kernel is tied to a specific radial similarity
ansatz. Replacing such objects wholesale would be a redesign of the development, not a
migration. What is genuinely replaceable divides into three classes.

\emph{Replaceable now, because the statement matches.} Compact-parameter integration is the
clearest case: \TauCeti{} proves exactly the regularity statement the project needs and
proves it for a general compact parameter space with a weight, which is strictly more general
than the project's version. Smooth functions with prescribed support are another: the project
constructs its own bump and cutoffs from $\exp(-1/x)$, while both \Mathlib{} and \TauCeti{}
already provide them.

\emph{Replaceable after a statement change.} The project defines Sobolev spaces of all
orders on a cylinder as a lattice of words to obtain completeness without a completion; a
library Sobolev space is a predicate on distributions. Moving to the library notion would
change the meaning of the objects the estimates are about, so this is a refactoring of the
Euler development rather than a substitution, and it should be done (if at all) only with the
estimates in view.

\emph{Not replaceable today.} The weighted band-indexed coefficient classes $W_\alpha$,
$M_\alpha$, $S_\alpha$ of the Navier--Stokes construction, the pulse envelopes and their
logarithmic edge weights, and the exponent ledger \eqref{eq:sigma} are bespoke objects of
this proof. A library would not want them.

\subsection{Outcome of the migration}
\label{sec:migration-results}

The dependency move itself is complete: the repository resolves, \Mathlib{}'s oleans are
fetched from the cache (9{,}057 files, about 4.8\,GB of build artifacts, under a minute on
a warm connection), and \TauCeti{}'s own modules build from its published cache (5{,}521
jobs, nine seconds). What remains is the repair of the project's own 644k lines against the
new API.

We worked through the failures reported by the compiler file by file. The sites we
repaired, and the shape of each repair, are the following.

\begin{center}\small
\begin{longtable}{@{}p{0.36\textwidth}p{0.58\textwidth}@{}}
\toprule
file & repair \\
\midrule
\file{R3/LpNormTools} & \lean{MemLp} projections; \lean{eLpNorm\_add\_le} lost its
  measurability arguments; \lean{eLpNorm\_one\_eq\_lintegral\_enorm} gained one; one lemma
  reproved by cases on measurability \\
\file{R3/PairedKernelBound} & the same, plus \lean{eLpNorm\_smul\_le\_mul\_eLpNorm}
  reordered (the scalar field comes first) and \lean{eLpNorm\_norm} needing its hypothesis \\
\file{R3/WeightedInterpolation} & \lean{Real.HolderTriple} construction; a measurability
  hypothesis added to the interpolation lemma and threaded to its call site;
  \lean{eLpNorm\_eq\_eLpNorm'} now needs its function named explicitly \\
\file{R3/PressureFunctionals}, \file{R3/H3Products} & \lean{MemLp} projections replaced by
  \lean{hf.aestronglyMeasurable} \\
\file{R3/ComparisonCutoffs} & \lean{norm\_iteratedFDeriv\_fderiv} restated with an explicit
  type; \lean{Finset.prod\_le\_prod} replaced by \lean{gcongr} \\
\file{AppendixHeatResults} & the same \lean{gcongr} replacement for a product of linear factors \\
\file{R3/ComparisonFiniteEnergy} & \lean{MemLp} projections (\lean{hf.1} to
  \lean{hf.aestronglyMeasurable}), including on a difference \lean{hf.sub\ hg} \\
\file{EdgeWeightJets}, \file{WeightedRadialPrimitive} & \lean{Finset.prod\_le\_one} needs
  \lean{MulLeftMono}, which $\mathbb{R}$ lacks; replaced by writing the product of equal
  factors as a power and using \lean{pow\_le\_one\textsubscript{0}} \\
\file{SmoothCovariance}, \file{SquaredPartition} & \lean{contDiff\_apply} and
  \lean{contDiff\_apply\_apply} now take the order implicitly and the family explicitly:
  \lean{contDiff\_apply\ \mathbb{R}\ (fun\ \_\ =>\ \mathbb{R})\ i} \\
\file{TauCetiBridge} (new) & imports \TauCeti{} and restates its parametric-integration
  theorems at project names; requires \lean{[CompleteSpace\ F]} on the codomain \\
\bottomrule
\end{longtable}
\end{center}

\paragraph{Status.} \textbf{The \file{NavierStokes} library builds.} Under the pinned
toolchain, \lean{lake build NavierStokes} completes with 19{,}645 jobs and exit code $0$ and no
diagnostics of any kind: all 816 modules elaborate. The two challenge modules of
\file{ComparatorChallenges} build as well (18{,}136 jobs), so of the repository's three
targets only \file{Euler} fails. Fifteen rounds of repair took it from 228
modules to the whole library, and no theorem statement was modified --- every change is a
proof, an explicit implicit argument, or, in exactly one place, a measurability hypothesis
added to a private helper lemma (\lean{lpNorm\_ofReal}) because the library lemma it calls now
demands one.

Two things made the difference, and both are worth recording for anyone else doing this.
First, \emph{read the goals, do not guess them}: inserting \lean{trace\_state} before each
failing line showed that the set a \lean{mono} expects is usually an intermediate one
manufactured by an enclosing \lean{apply}, not the theorem's own set. On that evidence the
eleven failures in \file{NominalProfile} fell in one pass, after two earlier attempts had
guessed wrong and made the file worse. Second, \emph{the two \lean{mono} lemmas do not name
their sets alike}:
\[
  \lean{IntegrableOn.mono\_set}\ (h : \lean{IntegrableOn}\ f\ t\ \mu)\ (hst : s \subseteq t)
  \quad\text{and}\quad
  \lean{ContDiffOn.mono}\ (h : \lean{ContDiffOn}\ f\ s)\ \{t\}\ (hst : t \subseteq s),
\]
so the target is \lean{s} in the first and \lean{t} in the second. Seven of the eleven sites
needed \lean{(s\ :=\ ...)} and four \lean{(t\ :=\ ...)}; applying one convention to all of them
does nothing and looks as though the argument had been ignored.

The \file{Euler} library is a different story and we do not dress it up. Its 1{,}839 modules
stand at 42 elaborated; \file{Euler/EulerProof} rejects 28 sites, five of which need a
measurability lemma the development never had (see below), three are \lean{case hf}
hypotheses, two are \lean{ContinuousAdd} instances that no longer resolve, and eighteen are
one-offs spread over 20{,}755 lines.

